% =====================================================================
% rapport.tex - Rapport academique TS6
% Projet IoT Edge Bridge - Noel Junior Yando Fotso - EPHEC 2025-2026
% Compilation : latexmk -xelatex -shell-escape rapport.tex
% =====================================================================
\documentclass[10pt,a4paper,french,twoside=false,open=any]{scrreprt}

% =====================================================================
% bridge-preamble.tex
% Preambule partage des livrables LaTeX du projet IoT Edge Bridge
% Auteur : Noel Junior Yando Fotso - EPHEC TS6 - 2025-2026
% Compilation : XeLaTeX + biber (latexmk -xelatex -shell-escape)
% =====================================================================

% --- Encodage et langue ----------------------------------------------
\usepackage[utf8]{inputenc}      % LuaTeX/XeTeX-friendly (ignore mais sans erreur)
\usepackage[T1]{fontenc}
\usepackage[french]{babel}
\usepackage{microtype}
\usepackage{csquotes}

% --- Fontes (XeLaTeX) ------------------------------------------------
% Detection conditionnelle : si XeLaTeX/LuaLaTeX, on essaie Fira Sans / Source Serif Pro,
% avec fallback automatique sur Latin Modern si la fonte est indisponible.
\usepackage{ifxetex,ifluatex}
\newif\ifxetexorluatex
\ifxetex\xetexorluatextrue\else
  \ifluatex\xetexorluatextrue\else\xetexorluatexfalse\fi
\fi

\ifxetexorluatex
  \usepackage{fontspec}
  % Source Serif Pro pour le rapport et le CDC ; Fira Sans pour le pitch.
  % Fallback silencieux sur Latin Modern si la fonte n'est pas installee.
  \IfFontExistsTF{Source Serif Pro}{%
    \setmainfont{Source Serif Pro}%
  }{%
    \IfFontExistsTF{Source Serif 4}{\setmainfont{Source Serif 4}}{}%
  }
  \IfFontExistsTF{Fira Sans}{%
    \setsansfont{Fira Sans}%
  }{}
  \IfFontExistsTF{Fira Mono}{%
    \setmonofont{Fira Mono}%
  }{}
\else
  % pdfLaTeX fallback : Latin Modern (par defaut, rien a faire).
\fi

% --- Graphiques et couleurs ------------------------------------------
\usepackage{graphicx}
\usepackage[table,dvipsnames,svgnames]{xcolor}

% Palette projet (charte)
\definecolor{bridgeblue}{HTML}{1F3A68}     % bleu profond, titres et cadres
\definecolor{bridgeteal}{HTML}{0FA4AF}     % teal d'accent
\definecolor{bridgegray}{HTML}{4A4A4A}     % gris texte secondaire
\definecolor{bridgealert}{HTML}{C62828}    % rouge alerte
\definecolor{bridgesoft}{HTML}{F2F6FB}     % bleu tres pale, fond doux
\definecolor{bridgeinfo}{HTML}{E3F4F4}     % teal tres pale, info
\definecolor{bridgewarn}{HTML}{FFF4E5}     % ambre tres pale, attention
\definecolor{bridgegreen}{HTML}{2E7D32}    % vert validation
\definecolor{bridgegreensoft}{HTML}{E8F5E9}% vert tres pale

% --- Liens hypertexte ------------------------------------------------
\usepackage[
  colorlinks=true,
  linkcolor=bridgeblue,
  citecolor=bridgeteal,
  urlcolor=bridgeteal,
  pdfborder={0 0 0}
]{hyperref}

% --- Tableaux et listes ----------------------------------------------
\usepackage{booktabs}
\usepackage{longtable}
\usepackage{tabularx}
\usepackage{array}
\usepackage{colortbl}
\usepackage{makecell}
\usepackage{enumitem}

% Style global de tableaux : zebrures bleu pale, regle haute teal, regles mid claires
\setlength{\arrayrulewidth}{0.4pt}
\renewcommand{\arraystretch}{1.25}
\arrayrulecolor{bridgeblue!40}

% Commande utilitaire pour activer les zebrures sur un tableau
% Usage : \begin{table} \bridgezebra ... \begin{tabular}...
\newcommand{\bridgezebra}{\rowcolors{2}{bridgesoft}{white}}

% Commande pour styliser une ligne d'en-tete (a placer dans la 1re ligne)
% Usage : \bridgeheader \textbf{Col1} & \textbf{Col2} \\ \midrule
\newcommand{\bridgeheader}{\rowcolor{bridgeblue!90}\color{white}\bfseries}

% Activation globale des zebrures par defaut sur tous les tableaux suivants.
% Pour desactiver localement : \rowcolors{1}{}{}
\rowcolors{2}{bridgesoft}{white}

% Style booktabs renforce avec couleur de regle
\let\oldtoprule\toprule
\let\oldmidrule\midrule
\let\oldbottomrule\bottomrule
\renewcommand{\toprule}{\arrayrulecolor{bridgeblue}\specialrule{0.7pt}{0pt}{0pt}\arrayrulecolor{bridgeblue!40}}
\renewcommand{\midrule}{\specialrule{0.4pt}{0pt}{0pt}}
\renewcommand{\bottomrule}{\arrayrulecolor{bridgeblue}\specialrule{0.7pt}{0pt}{0pt}\arrayrulecolor{bridgeblue!40}}

% Type de colonne avec en-tete bold blanc sur fond bleu
\newcolumntype{H}{>{\columncolor{bridgeblue}\color{white}\bfseries}l}

% Geometrie raisonnable (sera surchargee par les classes si besoin)
\usepackage{geometry}
\geometry{a4paper, margin=2.5cm, headheight=14pt}

% --- Boites informatives (tcolorbox) ---------------------------------
\usepackage[most]{tcolorbox}
\tcbset{
  bridgebox/.style={
    enhanced,
    colback=bridgesoft,
    colframe=bridgeblue,
    boxrule=0.4pt,
    arc=2pt,
    left=8pt,right=8pt,top=6pt,bottom=6pt,
    fonttitle=\bfseries\color{white},
    coltitle=white,
    colbacktitle=bridgeblue,
    title filled,
    breakable
  },
  bridgeinfobox/.style={
    bridgebox,
    colback=bridgeinfo,
    colframe=bridgeteal,
    colbacktitle=bridgeteal
  },
  bridgewarnbox/.style={
    bridgebox,
    colback=bridgewarn,
    colframe=orange!80!black,
    colbacktitle=orange!80!black
  },
  bridgekeybox/.style={
    bridgebox,
    colback=bridgegreensoft,
    colframe=bridgegreen,
    colbacktitle=bridgegreen
  }
}

% Macros raccourcies. title={#1} pour proteger les virgules dans le titre.
\newtcolorbox{bridgenote}[1][Note]{bridgebox, title={#1}}
\newtcolorbox{bridgeinfo}[1][Information]{bridgeinfobox, title={#1}}
\newtcolorbox{bridgewarn}[1][Attention]{bridgewarnbox, title={#1}}
\newtcolorbox{bridgekey}[1][A retenir]{bridgekeybox, title={#1}}

% --- Typographie KOMA-Script ----------------------------------------
% S'applique uniquement si on est en classe scrreprt/scrartcl/scrbook.
\makeatletter
\@ifclassloaded{scrreprt}{%
  \addtokomafont{disposition}{\color{bridgeblue}}%
  \addtokomafont{chapter}{\color{bridgeblue}\Huge\bfseries}%
  \addtokomafont{section}{\color{bridgeblue}}%
  \addtokomafont{subsection}{\color{bridgeblue}}%
  \addtokomafont{descriptionlabel}{\color{bridgeblue}\bfseries}%
  \KOMAoptions{numbers=noenddot, captions=tableheading, headings=optiontoheadandtoc}%
}{}
\@ifclassloaded{scrartcl}{%
  \addtokomafont{disposition}{\color{bridgeblue}}%
}{}
\makeatother

% --- Code et listings ------------------------------------------------
\usepackage{listings}

% Style generique
\lstdefinestyle{bridgebase}{%
  basicstyle=\ttfamily\small,
  keywordstyle=\color{bridgeblue}\bfseries,
  commentstyle=\color{bridgegray}\itshape,
  stringstyle=\color{bridgeteal},
  numberstyle=\tiny\color{bridgegray},
  numbers=left,
  numbersep=8pt,
  frame=single,
  rulecolor=\color{bridgegray!40},
  breaklines=true,
  showstringspaces=false,
  tabsize=2,
  captionpos=b
}

% Style JSON
\lstdefinelanguage{json}{%
  string=[s]{"}{"},
  stringstyle=\color{bridgeteal},
  comment=[l]{//},
  commentstyle=\color{bridgegray}\itshape,
  morestring=[b]",
  literate=
    *{0}{{{\color{bridgeblue}0}}}{1}
     {1}{{{\color{bridgeblue}1}}}{1}
     {2}{{{\color{bridgeblue}2}}}{1}
     {3}{{{\color{bridgeblue}3}}}{1}
     {4}{{{\color{bridgeblue}4}}}{1}
     {5}{{{\color{bridgeblue}5}}}{1}
     {6}{{{\color{bridgeblue}6}}}{1}
     {7}{{{\color{bridgeblue}7}}}{1}
     {8}{{{\color{bridgeblue}8}}}{1}
     {9}{{{\color{bridgeblue}9}}}{1}
     {:}{{{\color{bridgegray}{:}}}}{1}
     {,}{{{\color{bridgegray}{,}}}}{1}
     {\{}{{{\color{bridgeblue}{\{}}}}{1}
     {\}}{{{\color{bridgeblue}{\}}}}}{1}
     {[}{{{\color{bridgeblue}{[}}}}{1}
     {]}{{{\color{bridgeblue}{]}}}}{1},
}
\lstdefinestyle{bridgejson}{style=bridgebase, language=json}

% Style Python (utilise le langage builtin de listings)
\lstdefinestyle{bridgepython}{%
  style=bridgebase,
  language=Python,
  morekeywords={async,await,match,case}
}

\lstset{style=bridgebase}

% --- TikZ -------------------------------------------------------------
\usepackage{tikz}
\usetikzlibrary{shapes.geometric, arrows.meta, positioning, fit, backgrounds, calc}

% Styles partages pour les schemas du projet
\tikzset{
  bridgeblock/.style={%
    rectangle, rounded corners=2pt,
    draw=bridgeblue, thick,
    fill=bridgeblue!8,
    text=bridgeblue,
    align=center,
    minimum width=28mm, minimum height=10mm,
    inner sep=4pt,
    font=\small\sffamily
  },
  bridgeflow/.style={%
    -{Stealth[length=3mm,width=2mm]},
    draw=bridgeteal, thick
  },
  bridgealertstyle/.style={%
    rectangle, rounded corners=2pt,
    draw=bridgealert, thick,
    fill=bridgealert!8,
    text=bridgealert,
    align=center,
    minimum width=28mm, minimum height=10mm,
    font=\small\sffamily\bfseries
  },
  bridgenotetikz/.style={%
    rectangle, rounded corners=2pt,
    draw=bridgegray, dashed,
    fill=bridgegray!5,
    text=bridgegray,
    align=left,
    inner sep=4pt,
    font=\footnotesize\sffamily\itshape
  }
}

% Alias plus courts pour rester compatibles avec les schemas
\tikzset{bridgealert/.style={bridgealertstyle}}

% --- Bibliographie ---------------------------------------------------
\usepackage[
  backend=biber,
  style=ieee,
  sorting=nyt,
  maxbibnames=99,
  giveninits=true
]{biblatex}

% --- Sous-titre partage pour la page de garde ------------------------
% Les classes scrreprt et beamer ont leur propre logique de \subtitle ;
% on garde une commande explicite \bridgesubtitle pour compatibilite.
\providecommand{\bridgesubtitle}[1]{\gdef\@bridgesubtitle{#1}}
\providecommand{\thebridgesubtitle}{\@bridgesubtitle}

% Si la classe definit \subtitle (scrreprt KOMA, beamer), on l'utilise aussi.
\providecommand{\subtitle}[1]{\bridgesubtitle{#1}}

% --- Metadonnees PDF (renseignees a la compilation) ------------------
\AtBeginDocument{%
  \hypersetup{%
    pdfauthor={Noel Junior Yando Fotso},
    pdfsubject={IoT Edge Bridge - EPHEC TS6 2025-2026},
    pdfkeywords={IoT, FHIR, HL7, KMEHR, BIHR, MDR, RGPD, eSante, EPHEC}
  }
}

\addbibresource{references.bib}

% --- Mode compact pour le rapport (cible 15 pages) -------------------
% Marges serrees, parskip reduit, chapitres sans saut de page,
% titres compactes, espacement reduit autour des floats.
\geometry{a4paper, margin=1.6cm, headheight=14pt}

\setlength{\parskip}{0.4em plus 0.1em minus 0.05em}
\setlength{\parindent}{0pt}

% Reduit l'espace avant et apres les chapitres et sections
\RedeclareSectionCommand[
  beforeskip=-1.5em,
  afterskip=0.6em,
  font=\Large\bfseries\color{bridgeblue}
]{chapter}

\RedeclareSectionCommand[
  beforeskip=0.8em plus 0.1em minus 0.1em,
  afterskip=0.4em,
  font=\large\bfseries\color{bridgeblue}
]{section}

\RedeclareSectionCommand[
  beforeskip=0.6em plus 0.1em minus 0.1em,
  afterskip=0.3em,
  font=\normalsize\bfseries\color{bridgeblue}
]{subsection}

% Pas de saut de page du tout : les chapitres s'enchainent
\let\cleardoublepage\relax
\let\clearpage\relax
\renewcommand{\chapterpagestyle}{plain}

% Espace reduit autour des tableaux et figures
\setlength{\textfloatsep}{8pt plus 2pt minus 2pt}
\setlength{\intextsep}{6pt plus 2pt minus 2pt}
\setlength{\abovecaptionskip}{4pt}
\setlength{\belowcaptionskip}{4pt}

% Listes plus compactes
\setlist{itemsep=1pt, parsep=1pt, topsep=2pt, partopsep=0pt}

% Tableaux plus serres pour le mode compact
\renewcommand{\arraystretch}{1.05}

% Force la fonte des tableaux plus petite
\usepackage{etoolbox}
\AtBeginEnvironment{tabular}{\footnotesize}
\AtBeginEnvironment{tabularx}{\footnotesize}

% Bibliographie plus compacte
\renewcommand*{\bibfont}{\footnotesize}


% --- Metadonnees -----------------------------------------------------
\title{Passerelle d'Interop\'erabilit\'e Edge (IoT Bridge)\\
       pour l'int\'egration de dispositifs m\'edicaux Legacy au BIHR}
\subtitle{Rapport --- Technologies de la Sant\'e (TS6)}
\author{No\"el Junior Yando Fotso\\
        \small Encadrant : Matthieu OLIVIERS\\
        \small EPHEC Brussels}
\date{05/05/2026}

% =====================================================================
\begin{document}

% --- En-tete compact (one-column, pas de page de titre dediee) -------
\noindent
{\color{bridgeblue}\LARGE\bfseries Passerelle d'Interop\'erabilit\'e Edge (IoT Bridge)}\par
\vspace{2pt}
{\color{bridgeteal}\large pour l'int\'egration de dispositifs m\'edicaux Legacy au BIHR}\par
\vspace{6pt}
{\small Rapport, TS6 Technologies de la Sant\'e, EPHEC Brussels.\quad
 \textbf{Auteur} : No\"el Junior Yando Fotso \quad
 \textbf{Encadrant} : Matthieu OLIVIERS \quad
 \textbf{Date} : 05/05/2026}\par
\vspace{0.4em}\hrule height 0.6pt\vspace{0.6em}

% Pas de TOC pour atteindre 15 pages. Le rapport reste navigable via les
% titres de chapitre et les renvois croisés. La structure en 7 chapitres
% est annoncee dans l'introduction.

% --- Corps -----------------------------------------------------------
% =====================================================================
% Introduction generale du rapport
% =====================================================================
\chapter*{Introduction}
\addcontentsline{toc}{chapter}{Introduction}
\label{ch:introduction}

Le système de santé belge déploie le Belgian Integrated Health Record (BIHR), vue unique des données de santé du citoyen \cite{bihr_overview}. Le \emph{Plan d'action interfédéral eSanté 2025-2027} confie à l'INAMI son pilotage et au service eHealth la mise en {\oe}uvre technique \cite{plan_esante_2025_2027,ehealth_services_base}.

Un fossé technologique persiste. La majorité des hôpitaux belges échantillonnés n'atteignent pas le stade 4 de maturité numérique, qui suppose l'intégration automatique des dispositifs au DPI \cite{kce_362_dpi}. Trois causes : un parc Legacy de durée de vie 15 à 20 ans face à des standards qui évoluent tous les 3 à 5 ans ; le règlement MDR 2017/745 qui fait perdre le marquage CE à toute modification significative \cite{mdr_2017_745} ; l'absence de format universel d'export (RS-232, Bluetooth Classic, XML propriétaire, TCP brut coexistent).

Nous proposons un \emph{middleware Edge configurable et passif} qui absorbe l'hétérogénéité côté équipements et émet en HL7 FHIR R4 vers le DPI ou le BIHR \cite{hl7_fhir_r4}. La passerelle ne modifie ni le firmware ni le câblage : elle préserve le marquage CE et clarifie la chaîne de responsabilité. Configurable par profils JSON, elle ouvre la voie à une extensibilité communautaire.

\section*{Plan du rapport}

Le rapport suit la structure du cours TS6 :
\begin{enumerate}
  \item \emph{Introduction à l'e-health} : cadre belge, verrou Legacy, choix de FHIR.
  \item \emph{Applications mobiles et plateformes} : positionnement dans la pyramide d'architecture.
  \item \emph{Intégration des dispositifs} : c{\oe}ur technique. HL7 v2, FHIR, KMEHR, SNOMED-CT, LOINC, profils IHE PCD et DEC.
  \item \emph{Architecture technologique} : pile OSI, mTLS, services ETEE, perspectives SDN.
  \item \emph{Capteurs et dispositifs} : typologie et capture passive.
  \item \emph{Implémentation, adoption, évaluation} : déploiement, statut MDR, indicateurs.
  \item \emph{Éthique, sécurité et confidentialité} : RGPD, triade CIA, gouvernance, AIPD.
\end{enumerate}
La conclusion articule trois niveaux : technique, systémique, prospectif. Les annexes regroupent schémas, profils JSON et scénarios de démonstration.

% =====================================================================
% Chapitre 1 : Introduction a l'e-health et a la sante connectee
% =====================================================================
\chapter{Introduction à l'e-health et à la santé connectée}
\label{ch:ehealth}

Ce chapitre pose les définitions, situe la stratégie nationale eSanté, expose le verrou Legacy et justifie HL7 FHIR R4 comme standard pivot.

\section{Définition et domaines d'application}
\label{sec:ehealth-definition}

L'e-health désigne les technologies numériques médicales (apps mobiles, logiciels cliniques, capteurs, wearables) qui soutiennent diagnostic, monitoring, traitement et prévention. Le KCE caractérise ce périmètre dans son rapport sur le DPI hospitalier \cite{kce_362_dpi}. L'OMS structure le champ par sa \emph{Classification of Digital Health Interventions v2.0} \cite{oms_dhi_2023} : catégorie 1 \emph{pour les patients} (auto-monitoring), catégorie 2 \emph{pour les prestataires} (téléconsultation, transmission diagnostique, aide à la décision), catégorie 3 \emph{pour les gestionnaires} (logistique, registres).

\begin{bridgeinfo}[Positionnement de notre projet]
La passerelle relève sans ambiguïté de la catégorie 2 de l'OMS, sous-classes \emph{transmission of diagnostic data} et \emph{remote patient monitoring}. Elle s'adresse aux prestataires (cliniciens, techniciens biomédicaux, équipes IT), pas au patient. Elle ne porte pas d'interface clinique : son rôle est d'alimenter le DPI et le BIHR à partir d'équipements existants.
\end{bridgeinfo}

\section{Stratégie e-santé belge et plan eSanté 2025-2027}
\label{sec:ehealth-belge}

Le \emph{Plan eSanté 2025-2027} érige le BIHR en colonne vertébrale de la santé numérique belge \cite{plan_esante_2025_2027}. L'INAMI assume la responsabilité institutionnelle. Le service eHealth.fgov.be opère les briques transverses : authentification eID, services ETEE, eHealthBox, MyCareNet, horodatage \cite{ehealth_services_base}. Le BIHR vise une vue unique accessible aux soignants autorisés via relation thérapeutique vérifiée \cite{bihr_overview}. MetaHub fédère les hubs régionaux (RSW, ABRUMET, CoZo) et expose un index national des documents de santé \cite{metahub_overview}. La passerelle alimente le DPI hospitalier, qui publie ensuite via les hubs et MetaHub vers le BIHR.

\section{Problématique du matériel en exploitation (Legacy)}
\label{sec:ehealth-legacy}

Le parc Legacy est un goulet structurel. Un dispositif critique a une durée de vie de 15 à 20 ans, alors que les standards évoluent tous les 3 à 5 ans. Un moniteur acquis en 2008 a été conçu à l'époque de RS-232 et Bluetooth Classic ; la stratégie actuelle s'appuie sur FHIR R4 et OAuth 2.0. Le règlement MDR 2017/745 prévoit qu'une modification significative de la conception ou de la destination entraîne une nouvelle évaluation et la perte du marquage CE \cite{mdr_2017_745}. Ajouter une connectivité native suppose collaboration du fabricant, révision du dossier technique et nouvelle certification : coût et délai prohibitifs pour des centaines d'équipements.

\begin{bridgewarn}[Le paradoxe du \emph{Locked Config}]
Beaucoup d'équipements actifs s'appuient encore sur RS-232, RS-485 ou Bluetooth Classic, standards de leur époque (2005 à 2015). Ils tournent souvent sur des OS embarqués figés (Windows XP Embedded, Windows 7 Embedded), gelés par le fabricant pour garantir la sécurité du patient et préserver la validation logicielle. Mettre à jour l'OS reviendrait à modifier le dispositif, donc à perdre la certification.
\end{bridgewarn}

\section{Constat de maturité numérique en Belgique}
\label{sec:ehealth-maturite}

Selon les modèles de maturité numérique hospitalière comme HIMSS EMRAM, peu d'établissements belges atteignent à ce jour le stade d'intégration automatique des dispositifs médicaux dans le parcours de soin. Le KCE Report 362 \cite{kce_362_dpi} confirme cette inertie sans fournir de chiffre consolidé. Le déficit est multifactoriel : MDR ralentit les évolutions, l'absence de format universel impose des intégrations point à point, le coût d'un parc neuf est non finançable à court terme, les compétences d'intégration biomédicale sont rares.

\section{Avantages attendus de l'intégration automatisée}
\label{sec:avantages-attendus}

L'intégration automatique des dispositifs médicaux apporte quatre bénéfices documentés par l'OMS \cite{oms_dhi_2023}.

\begin{table}[h]
  \centering
  \caption{Bénéfices attendus.}
  \label{tab:benefices}
  \bridgezebra
  \begin{tabular}{p{4.0cm} p{10.5cm}}
    \toprule
    \bridgeheader Bénéfice & Effet attendu \\
    \midrule
    Qualité des soins & Dossier patient complet, traçabilité des mesures \\
    Gain de temps soignant & Suppression de la saisie manuelle au lit \\
    Médecine prédictive & Données denses exploitables par modèles d'aide à la décision \\
    Réduction des coûts & Évite le remplacement du parc Legacy, baisse des erreurs de saisie \\
    \bottomrule
  \end{tabular}
\end{table}

\section{Choix de HL7 FHIR R4 comme standard pivot}
\label{sec:ehealth-fhir}

HL7 FHIR R4 (HL7 International, 2019) est le standard retenu \cite{hl7_fhir_r4}. Conçu pour l'API REST, il expose un modèle granulaire par ressources (\texttt{Patient}, \texttt{Device}, \texttt{Observation}, \texttt{DiagnosticReport}) et utilise JSON nativement. Cette granularité contraste avec les messages monolithiques de HL7 v2 et facilite l'agrégation différenciée dans le BIHR. FHIR R4 sera le socle du futur \emph{European Health Data Space} (EHDS), s'inscrit dans la continuité de HL7 v2 et CDA, et accepte des profils nationaux. La passerelle adopte FHIR comme sortie principale, avec option HL7 v2 (DPI Legacy) et KMEHR (hubs régionaux).

\section{Notre proposition : un middleware d'intégration}
\label{sec:ehealth-proposition}

La passerelle IoT Edge Bridge est un middleware Python configurable, déployé en bord de réseau. Elle absorbe quatre canaux d'entrée : TCP brut (moniteurs IP), BLE (dispositifs portables), fichiers XML (spiromètres, ECG), trames série RS-232 (pompes, ventilateurs). Elle émet en FHIR R4 vers HAPI ou le DPI, avec option HL7 v2 over MLLP. La capture est passive : ni firmware ni câblage modifiés. La passerelle contribue indirectement à l'empowerment du patient : les données vitales captées et publiées au DPI alimentent le portail MaSanté.be, qui rend visible au patient ses propres mesures.

\begin{bridgekey}[Statut de la passerelle]
Le middleware est un \textbf{outil d'intégration IT}, pas un dispositif médical. Il consomme des sorties existantes et les normalise. Le fabricant reste responsable du dispositif, l'éditeur de la passerelle est responsable du mapping, l'hôpital est responsable du traitement. La passerelle est régie par le RGPD et l'ISO 27001, sans procédure MDR de classe IIa.
\end{bridgekey}

\begin{figure}[h]
  \centering
  \resizebox{\linewidth}{!}{% =====================================================================
% figures/fig-chaos-ordre.tex
% Vision narrative : du chaos des dispositifs au continuum BIHR.
% Trois zones : chaos (rouge), bridge (bleu, 4 couches avec description),
% ordre (vert).
% =====================================================================
\begin{tikzpicture}[
  font=\sffamily,
  every node/.style={font=\small\sffamily},
  >={Stealth[length=2.5mm]},
  thick,
  device/.style={
    rectangle, rounded corners=2pt,
    draw=bridgealert, fill=bridgealert!10, text=bridgealert,
    minimum width=2.4cm, minimum height=0.7cm,
    inner sep=2pt, font=\small\sffamily
  },
  proto/.style={
    font=\scriptsize\itshape, color=bridgegray,
    inner sep=1pt
  },
  layer/.style={
    rectangle, rounded corners=2pt,
    draw=bridgeblue, fill=bridgeblue!10, text=bridgeblue,
    minimum width=2.0cm, minimum height=0.7cm,
    inner sep=2pt, font=\small\bfseries\sffamily
  },
  layerdesc/.style={
    text=bridgegray, align=left,
    inner sep=1pt,
    font=\scriptsize\sffamily
  },
  target/.style={
    rectangle, rounded corners=2pt,
    draw=bridgegreen, fill=bridgegreen!10, text=bridgegreen,
    minimum width=3.2cm, minimum height=0.7cm,
    inner sep=2pt, font=\small\sffamily
  },
  zonebox/.style={rounded corners=4pt, thick, fill opacity=1},
  flowarrow/.style={-{Stealth[length=3mm]}, very thick}
]

% =====================================================================
% Zone 1 : Chaos (gauche) - dispositifs alignes verticalement
% =====================================================================
\node[zonebox, draw=bridgealert!70, dashed, fill=bridgealert!4,
      minimum width=5.0cm, minimum height=6.4cm,
      anchor=north west, label={[anchor=north, font=\small\bfseries,
                                 color=bridgealert, yshift=-4pt]north:Chaos cote sources}]
      (chaoszone) at (0, 6.2) {};

\node[device] (d1) at (1.4, 4.8) {Moniteur};
\node[device] (d2) at (1.4, 4.05) {Glucometre};
\node[device] (d3) at (1.4, 3.30) {Pompe};
\node[device] (d4) at (1.4, 2.55) {Logiciel labo};
\node[device] (d5) at (1.4, 1.80) {DMS Legacy};
\node[device] (d6) at (1.4, 1.05) {Export Excel};

\node[proto, anchor=west] at ([xshift=2pt]d1.east) {RS-232};
\node[proto, anchor=west] at ([xshift=2pt]d2.east) {BLE/JSON};
\node[proto, anchor=west] at ([xshift=2pt]d3.east) {MEDIBUS};
\node[proto, anchor=west] at ([xshift=2pt]d4.east) {HL7v2};
\node[proto, anchor=west] at ([xshift=2pt]d5.east) {XML maison};
\node[proto, anchor=west] at ([xshift=2pt]d6.east) {CSV/XLSX};

% =====================================================================
% Zone 2 : Bridge (centre) - 4 couches AVEC description courte a droite
% =====================================================================
\node[zonebox, draw=bridgeblue, fill=bridgeblue!4,
      minimum width=5.4cm, minimum height=6.4cm,
      anchor=north west, label={[anchor=north, font=\small\bfseries,
                                 color=bridgeblue, yshift=-4pt]north:IoT Edge Bridge}]
      (bridgezone) at (6.0, 6.2) {};

% Coordonnees X : couche a x=7.0, description a x=8.6
\def\xL{7.0}
\def\xD{8.6}

\node[layer, anchor=center] (adapter) at (\xL, 4.95) {Adapter};
\node[layerdesc, anchor=west] at (\xD, 4.95)
  {Capte le canal :\\port serie, MQTT,\\TCP, fichier};

\node[layer, anchor=center] (parser) at (\xL, 3.85) {Parser};
\node[layerdesc, anchor=west] at (\xD, 3.85)
  {Decode la trame :\\HL7v2, JSON,\\MEDIBUS, CSV};

\node[layer, anchor=center] (mapper) at (\xL, 2.75) {Mapper};
\node[layerdesc, anchor=west] at (\xD, 2.75)
  {Codifie en LOINC,\\ecrit le Bundle\\FHIR R4};

\node[layer, anchor=center] (transp) at (\xL, 1.65) {Transport};
\node[layerdesc, anchor=west] at (\xD, 1.65)
  {Publie en mTLS\\avec retry et\\fallback SQLite};

\draw[bridgeflow] (adapter.south) -- (parser.north);
\draw[bridgeflow] (parser.south)  -- (mapper.north);
\draw[bridgeflow] (mapper.south)  -- (transp.north);

\node[font=\scriptsize\itshape, color=bridgegray, anchor=north]
  at (8.2, 1.05) {Profils JSON + plugins};

% =====================================================================
% Zone 3 : Ordre (droite) - destinataires alignes verticalement
% =====================================================================
\node[zonebox, draw=bridgegreen!70, fill=bridgegreen!4,
      minimum width=4.6cm, minimum height=6.4cm,
      anchor=north west, label={[anchor=north, font=\small\bfseries,
                                 color=bridgegreen, yshift=-4pt]north:Ordre cote plateforme}]
      (ordrezone) at (12.0, 6.2) {};

\node[target] (fhir) at (14.3, 4.8) {HAPI FHIR R4};
\node[target] (dpi)  at (14.3, 3.9) {DPI hospitalier};
\node[target] (bihr) at (14.3, 3.0) {MetaHub / BIHR};
\node[target] (cli)  at (14.3, 2.1) {Clinicien};

\draw[bridgeflow] (fhir.south) -- (dpi.north);
\draw[bridgeflow] (dpi.south)  -- (bihr.north);
\draw[bridgeflow] (bihr.south) -- (cli.north);

% =====================================================================
% Fleches inter-zones, parfaitement horizontales au niveau du milieu
% =====================================================================
\draw[flowarrow, draw=bridgegray]
  (chaoszone.east |- 0,3.5) -- (bridgezone.west |- 0,3.5)
  node[midway, above, font=\scriptsize\itshape, color=bridgegray] {capture};

\draw[flowarrow, draw=bridgegray]
  (bridgezone.east |- 0,3.5) -- (ordrezone.west |- 0,3.5)
  node[midway, above, font=\scriptsize\itshape, color=bridgegray] {publication};

\end{tikzpicture}
}
  \caption{Vision narrative du projet : du chaos des dispositifs Legacy vers un continuum BIHR uniformisé.}
  \label{fig:rapport-vision}
\end{figure}

% =====================================================================
% Chapitre 2 : Applications mobiles et plateformes de sante connectee
% =====================================================================
\chapter{Applications mobiles et plateformes de santé connectée}
\label{ch:architecture-applicative}

Ce chapitre positionne la passerelle dans la pyramide d'architecture du SI hospitalier, identifie le DPI comme cible, situe le traitement sur le continuum DIKW et précise le rôle du dashboard.

\section{Processus métier global}
\label{sec:processus-metier}

Le diagramme suivant résume le processus métier que la passerelle automatise. Une source (dispositif IoT ou logiciel Legacy) émet des données, le bridge les capture, les normalise et les valide, le DPI les persiste, et le soignant les consulte ensuite via son interface habituelle.

\begin{figure}[h]
  \centering
  \resizebox{\linewidth}{!}{% =====================================================================
% figures/fig-processus-metier.tex
% Diagramme de processus metier BPMN simplifie : source -> bridge -> DPI -> soignant.
% Swimlane horizontale a 4 lignes.
% =====================================================================
\begin{tikzpicture}[
  font=\sffamily,
  every node/.style={font=\footnotesize\sffamily},
  >={Stealth[length=2.5mm]},
  thick,
  source/.style={
    rectangle, rounded corners=2pt,
    draw=bridgealert, fill=bridgealert!10, text=bridgealert,
    minimum width=2.4cm, minimum height=0.8cm, align=center,
    inner sep=2pt, font=\footnotesize\sffamily
  },
  bridge/.style={
    rectangle, rounded corners=2pt,
    draw=bridgeblue, fill=bridgeblue!10, text=bridgeblue,
    minimum width=2.4cm, minimum height=0.8cm, align=center,
    inner sep=2pt, font=\footnotesize\sffamily
  },
  dpi/.style={
    rectangle, rounded corners=2pt,
    draw=bridgeteal, fill=bridgeteal!12, text=bridgeteal,
    minimum width=2.4cm, minimum height=0.8cm, align=center,
    inner sep=2pt, font=\footnotesize\sffamily
  },
  user/.style={
    rectangle, rounded corners=2pt,
    draw=bridgegreen, fill=bridgegreen!12, text=bridgegreen,
    minimum width=2.4cm, minimum height=0.8cm, align=center,
    inner sep=2pt, font=\footnotesize\sffamily
  },
  decision/.style={
    diamond, draw=bridgeblue, fill=bridgeblue!8, text=bridgeblue,
    minimum width=1.4cm, minimum height=1.0cm, inner sep=1pt,
    font=\scriptsize\sffamily, aspect=1.6
  },
  flow/.style={-{Stealth[length=3mm]}, thick, draw=bridgeblue!80},
  errflow/.style={-{Stealth[length=3mm]}, thick, dashed, draw=bridgealert}
]

% --- Lignes (swim lanes) ---------------------------------------------
\foreach \y/\name in {3.6/Source, 2.4/Bridge, 1.2/DPI, 0/Soignant}
  \node[anchor=east, font=\footnotesize\bfseries\sffamily, color=bridgegray]
    at (-0.2, \y) {\name};

% Lignes de separation
\foreach \y in {3.0, 1.8, 0.6}
  \draw[bridgegray!30, dashed] (-0.1,\y) -- (16.5,\y);

% --- Bandeau Source --------------------------------------------------
\node[source] (s1) at (1.2, 3.6) {Dispositif IoT\\ou logiciel Legacy};
\node[font=\scriptsize\itshape, color=bridgegray]
  at (s1.south)[anchor=north, yshift=-1pt] {emet trame};

% --- Bandeau Bridge --------------------------------------------------
\node[bridge] (b1) at (4.0, 2.4) {Capture\\Adapter};
\node[bridge] (b2) at (6.4, 2.4) {Parsing\\Mapping};
\node[decision] (d1) at (8.8, 2.4) {Valide ?};
\node[bridge] (b3) at (11.2, 2.4) {Transport\\mTLS};

% Quarantaine en cas d'erreur
\node[bridge, fill=bridgealert!10, draw=bridgealert, text=bridgealert]
  (q) at (8.8, 3.6) {Quarantaine\\+ alerte};

% --- Bandeau DPI -----------------------------------------------------
\node[dpi] (dpi1) at (13.6, 1.2) {DPI hospitalier\\persiste FHIR};

% --- Bandeau Soignant ------------------------------------------------
\node[user] (u1) at (15.6, 0) {Consultation\\au DPI};

% --- Fleches du flux nominal -----------------------------------------
\draw[flow] (s1.south) -- (b1.north -| s1.south);
\draw[flow] (s1.south) |- (b1.west);
\draw[flow] (b1.east) -- (b2.west);
\draw[flow] (b2.east) -- (d1.west);
\draw[flow] (d1.east) -- (b3.west) node[midway, above, font=\scriptsize] {oui};
\draw[flow] (b3.south) |- (dpi1.west);
\draw[flow] (dpi1.south) |- (u1.west);

% --- Fleche d'erreur -------------------------------------------------
\draw[errflow] (d1.north) -- (q.south) node[midway, right, font=\scriptsize, color=bridgealert] {non};

\end{tikzpicture}
}
  \caption{Processus métier : de la source Legacy à la consultation clinique.}
  \label{fig:rapport-processus-metier}
\end{figure}

\section{Pyramide d'architecture des systèmes d'information}
\label{sec:pyramide}

\begin{table}[h]
  \centering
  \caption{Pyramide d'architecture appliquée à la passerelle.}
  \label{tab:pyramide}
  \bridgezebra
  \begin{tabular}{p{2.4cm} p{5.6cm} p{6.0cm}}
    \toprule
    \bridgeheader Niveau & Préoccupations & Position de la passerelle \\
    \midrule
    Métier & Vision, valeur clinique, finalité de soin & Améliore la continuité et la complétude des données du patient \\
    Fonctionnel & Processus métiers (urgences, hospitalisation, bloc) & Automatise la capture des observations vers le DPI \\
    Applicatif & Logiciels DPI, PACS, LIS, RIS & Middleware entre les dispositifs et le DPI hospitalier \\
    Technique & Serveurs, réseau, dispositifs physiques & Gateway edge Python exécutée sur poste local hospitalier \\
    \bottomrule
  \end{tabular}
\end{table}

La passerelle agit principalement au niveau applicatif. Elle s'appuie sur le niveau technique (port série, BLE, TCP, fichier) et sert le niveau fonctionnel en supprimant la saisie manuelle. Sa contribution au métier est indirecte mais réelle : un dossier plus complet améliore la décision clinique. Dans la chaîne canonique \emph{capteur, microcontrôleur, communication, application, actionneur}, notre passerelle se substitue au maillon \emph{application} (cf. figure~\ref{fig:rapport-chaine-capteur} au chapitre~\ref{ch:reseau}).

\section{Le DPI hospitalier : cible de notre passerelle}
\label{sec:dpi}

Le DPI est le point de convergence de toutes les données cliniques. Le KCE le structure autour de huit fonctionnalités : collecte, partage, sécurité d'accès, interopérabilité externe, aide à la décision, prescriptions, suivi longitudinal, portail patient \cite{kce_362_dpi}. Le principe \emph{only once} (saisie unique à la source) gouverne sa conception. Trois éditeurs structurent le marché belge : Omnipro, H+, Xcare. Tous exposent désormais une API FHIR locale, conforme au Plan eSanté 2025-2027 \cite{plan_esante_2025_2027}. La passerelle alimente cette API : flux dispositifs transformés en \texttt{Observation} liées au \texttt{Patient} et au \texttt{Device}. En démonstration, un serveur HAPI FHIR local simule cette interface.

\section{Continuum d'intelligence des données}
\label{sec:dikw}

Le modèle DIKW (Données, Information, Connaissance, Sagesse) décrit la chaîne d'une mesure brute vers une décision. La passerelle opère sur les deux premiers maillons.

\begin{table}[h]
  \centering
  \caption{Continuum DIKW illustré sur la passerelle.}
  \label{tab:dikw}
  \bridgezebra
  \begin{tabular}{p{3.8cm} p{10.0cm}}
    \toprule
    \bridgeheader Niveau DIKW & Illustration projet \\
    \midrule
    Donnée brute & Trame RS-232 ASCII \texttt{0x32 0x35 0x6D ...} émise par le moniteur Philips \\
    Information & SpO2 = 96\% à 14:23 sur le lit 12, dispositif Philips MX450 \\
    Connaissance & SpO2 stable sur 24 heures, dans la plage de plausibilité 50 à 100\% \\
    Décision clinique & Pas d'alerte, suivi standard. Hors-scope passerelle, opéré par le clinicien sur le DPI \\
    \bottomrule
  \end{tabular}
\end{table}

Cette frontière est volontaire. La passerelle ne juge pas la gravité d'une mesure et n'émet pas d'alerte. Elle sépare la responsabilité du fabricant (justesse de la mesure), celle de l'éditeur (justesse du mapping) et celle du clinicien (interprétation).

\section{Algorithmes et paradigmes mobilisés}
\label{sec:algorithmes}

La passerelle est déterministe. Le parsing applique un automate à états finis sur les trames série, le mapping applique des règles déclaratives lues depuis des profils JSON, le transport applique retry et buffering. Aucun modèle d'apprentissage n'est embarqué dans la démonstration. Cette simplicité facilite validation, traçabilité et reproductibilité. L'IA en bord (détection d'anomalies, pré-flagging) figure en roadmap F2 du cahier des charges.

\section{Interface de supervision OT}
\label{sec:dashboard}

La passerelle expose un dashboard FastAPI + HTMX destiné aux techniciens biomédicaux et aux équipes IT. Il remplit deux fonctions : configuration des profils JSON par équipement (canal, langue, mapping de sortie) et supervision OT (statut couleur, horodatage de dernière communication, nombre de messages, contexte technique en cas d'anomalie). Les notifications mail respectent un format strict (identifiant équipement, type d'erreur, horodatage, contexte technique) sans aucun paramètre clinique.

\begin{bridgekey}[Principe de minimisation]
Le dashboard regarde les flux, pas le contenu clinique. Les valeurs vitales transitent vers le DPI sans s'arrêter dans le tableau de bord. Les techniciens biomédicaux et OT voient \emph{que} les flux circulent, jamais \emph{ce qui} circule. Cette règle préserve la confidentialité patient et clarifie la mission OT du middleware.
\end{bridgekey}

\section{Écosystème applicatif belge environnant}
\label{sec:ecosysteme-belge}

La passerelle alimente le DPI, qui s'inscrit dans un écosystème applicatif belge plus large. Cinq briques transversales structurent cet écosystème.

\begin{table}[h]
  \centering
  \caption{Briques applicatives belges et leur position vis-à-vis de la passerelle.}
  \label{tab:ecosysteme-be}
  \bridgezebra
  \begin{tabular}{p{3.0cm} p{11.5cm}}
    \toprule
    \bridgeheader Brique & Rôle et lien avec la passerelle \\
    \midrule
    Bingli & Chatbot d'anamnèse en amont du DPI, indépendant du flux passerelle \\
    Recip-e & E-prescription, alimentée par le DPI, hors-scope passerelle \\
    eHealthBox & Boîte chiffrée pour échanges inter-praticiens, alimentée par le DPI \\
    Sumehr & Résumé patient partageable via le RSW, alimenté par le DPI \\
    MyCareNet & Plateforme INAMI pour la facturation, hors-scope passerelle \\
    \bottomrule
  \end{tabular}
\end{table}

La passerelle se positionne en amont du DPI : elle alimente la couche transactionnelle. Les briques transversales s'appuient ensuite sur la complétude du DPI pour leurs propres traitements.

% =====================================================================
% Chapitre 3 : Integration des dispositifs connectes (CŒUR du rapport)
% =====================================================================
\chapter{Intégration des dispositifs connectés dans les systèmes de santé}
\label{ch:standards}

Comment traduire les sorties propriétaires des dispositifs Legacy en standards acceptés par le BIHR ? Ce chapitre est le c{\oe}ur du rapport : trois couches d'interopérabilité, standards mobilisés (HL7 v2, FHIR R4, IHE PCD, LOINC, KMEHR, DICOM, SNOMED-CT), architecture de mapping et insertion dans le paysage belge.

\section{Les trois couches d'interopérabilité}
\label{sec:interop-3couches}

\begin{table}[h]
  \centering
  \caption{Les trois couches d'interopérabilité.}
  \label{tab:3couches}
  \bridgezebra
  \begin{tabular}{p{2.8cm} p{6.2cm} p{5.0cm}}
    \toprule
    \bridgeheader Couche & Question résolue & Standards \\
    \midrule
    Technique & Comment les bits passent-ils d'un système à l'autre ? & RS-232, BLE, MLLP, HTTPS \\
    Syntaxique & Comment les bits sont-ils structurés en messages ? & HL7 v2, FHIR R4, KMEHR, DICOM \\
    Sémantique & Quel sens partagé donner aux mêmes mots ? & LOINC, SNOMED-CT, ICD-10 \\
    \bottomrule
  \end{tabular}
\end{table}

La passerelle traverse les trois couches : l'Adapter résout la couche technique, le Parser et le Mapper résolvent les couches syntaxique et sémantique. Cette répartition reprend la philosophie IHE \cite{ihe_pcd_tf}.

\section{Cinq défis structurels de l'interopérabilité}
\label{sec:defis-interop}

Au-delà des trois couches, le cours TS6 identifie cinq défis structurels que toute solution d'interopérabilité doit relever.

\begin{table}[h]
  \centering
  \caption{Cinq défis de l'interopérabilité et leur traitement par la passerelle.}
  \label{tab:5defis}
  \bridgezebra
  \begin{tabular}{p{4.5cm} p{10.0cm}}
    \toprule
    \bridgeheader Défi & Réponse de la passerelle \\
    \midrule
    Compatibilité technique entre fabricants & Adapters multi-protocoles, profils JSON par dispositif \\
    Échange et interprétation cohérents & Mapper FHIR R4 + codification LOINC \\
    Coordination des soins & Alimentation du DPI, base de la coordination clinique \\
    Réduction des erreurs médicales & Suppression de la saisie manuelle, validation FHIR avant émission \\
    Vue holistique du patient & Agrégation multi-sources dans le DPI puis le BIHR \\
    \bottomrule
  \end{tabular}
\end{table}

\section{HL7 v2 : le standard hospitalier historique}
\label{sec:hl7v2}

HL7 v2 (1989) reste le standard de messagerie clinique le plus déployé \cite{hl7_v2}. Il opère à la couche application OSI. Ses messages sont des chaînes ASCII pipe-delimited en segments : MSH (en-tête), PID (patient), OBR (requête), OBX (résultat). Variantes courantes : ORU\textasciicircum{}R01 (résultats), ADT (admissions), ORM (commandes). En démonstration, le simulateur Philips IntelliVue émet ORU\textasciicircum{}R01 over MLLP sur le port 2575. L'Adapter ouvre une socket TCP, reconnaît les délimiteurs MLLP (\texttt{0x0B}, \texttt{0x1C 0x0D}) et transmet au Parser, qui extrait les segments OBX et passe au Mapper. La passerelle convertit en FHIR R4, format pivot.

\section{HL7 FHIR R4 : la modernité REST}
\label{sec:fhir}

HL7 FHIR R4 (2019) raisonne en ressources atomiques accessibles par API REST \cite{hl7_fhir_r4}. Ressources mobilisées : \texttt{Patient}, \texttt{Device}, \texttt{Observation}, \texttt{DeviceMetric}, \texttt{Bundle}. Format natif JSON. Chaque ressource peut être profilée et validée. La passerelle émet des Bundle de type \texttt{transaction} contenant \texttt{Patient}, \texttt{Device} et \texttt{Observation}. Le serveur cible (HAPI FHIR en démo, Omnipro ou H+ en production) persiste les ressources et renvoie les identifiants permanents. Un exemple complet figure en annexe B du cahier des charges.

\section{IHE PCD : profil d'intégration des dispositifs}
\label{sec:ihe-pcd}

IHE orchestre l'usage coordonné de DICOM et HL7 par des profils d'intégration. Le domaine PCD (Patient Care Device) cible l'intégration des dispositifs au SI \cite{ihe_pcd_tf}. Le profil PCD-01, \emph{Communicate Patient Care Device Data}, formalise la transmission des données vitales. Notre Adapter joue le rôle de \emph{Device Observation Reporter} en transposant PCD-01 dans une implémentation Python configurable.

\section{Vocabulaires sémantiques : LOINC et SNOMED-CT}
\label{sec:loinc-snomed}

LOINC code chaque grandeur biologique ou clinique mesurable. SNOMED-CT couvre un périmètre plus large incluant diagnostics, procédures et contextes \cite{hl7_fhir_r4}. La passerelle utilise LOINC pour codifier les observations. SNOMED-CT figure en perspective.

\begin{table}[h]
  \centering
  \caption{Mesures principales et codes LOINC associés.}
  \label{tab:loinc}
  \bridgezebra
  \begin{tabular}{p{4.2cm} p{2.6cm} p{6.5cm}}
    \toprule
    \bridgeheader Mesure & Code LOINC & Profil source \\
    \midrule
    Saturation en oxygène (SpO2) & 59408-5 & Moniteur Philips IntelliVue \\
    Fréquence cardiaque & 8867-4 & Moniteur Philips IntelliVue \\
    Pression artérielle systolique & 8480-6 & Moniteur Philips IntelliVue \\
    Pression artérielle diastolique & 8462-4 & Moniteur Philips IntelliVue \\
    Glycémie capillaire & 2339-0 & Glucomètre Masimo via BLE \\
    Volume expiratoire forcé (VEMS) & 19868-9 & Spiromètre Vyaire MicroLab \\
    Débit perfusion & 96821-9 & Pompe B. Braun Infusomat \\
    \bottomrule
  \end{tabular}
\end{table}

\section{KMEHR : le standard belge}
\label{sec:kmehr}

KMEHR (Kind Messages for Electronic Healthcare Record) est le format pivot des échanges électroniques de santé en Belgique. Format XML structuré en transactions thématiques (résultats labo, lettres de sortie, prescriptions) \cite{kmehr_spec}. Format attendu par eHealthBox et les hubs MetaHub : RSW, ABRUMET, CoZo \cite{ehealth_services_base}. La passerelle ne produit pas de KMEHR en démo ; une couche de sortie KMEHR figure en roadmap.

\section{DICOM : pour mémoire}
\label{sec:dicom}

DICOM est le standard universel de l'imagerie médicale, couvrant format des images et protocoles d'échange entre modalités, PACS et stations \cite{dicom_standard}. La passerelle ne traite pas d'images en démo ; DICOM figure en roadmap pour le point-of-care imaging (échographie portable, ECG ambulatoire avec tracé).

\section{Architecture de mapping de la passerelle}
\label{sec:archi-mapping}

Pipeline à quatre couches : Adapter (flux physiques), Parser (structures propriétaires vers objets normalisés), Mapper (profil JSON, codification LOINC), Transport (sérialisation FHIR R4 et envoi).

\begin{figure}[h]
  \centering
  \resizebox{\linewidth}{!}{% =====================================================================
% figures/fig-pipeline-4couches.tex
% Pipeline fonctionnel : Adapter -> Parser -> Mapper -> Transport -> HAPI
% Inputs varies a gauche, sortie FHIR a droite. Bandeau plugins au-dessus.
% =====================================================================
\begin{tikzpicture}[
  font=\sffamily,
  every node/.style={font=\small\sffamily},
  >={Stealth[length=2.5mm]},
  thick,
  layer/.style={
    rectangle, rounded corners=3pt,
    draw=bridgeblue, fill=bridgeblue!10, text=bridgeblue,
    minimum width=2.4cm, minimum height=4.2cm,
    align=center,
    font=\small\bfseries\sffamily
  },
  altlayer/.style={
    rectangle, rounded corners=3pt,
    draw=bridgeteal, fill=bridgeteal!10, text=bridgeteal,
    minimum width=2.4cm, minimum height=4.2cm,
    align=center,
    font=\small\bfseries\sffamily
  },
  outbox/.style={
    rectangle, rounded corners=3pt,
    draw=bridgegreen, fill=bridgegreen!15, text=bridgegreen,
    minimum width=2.0cm, minimum height=1.0cm,
    align=center,
    font=\small\bfseries\sffamily
  },
  inputlbl/.style={
    anchor=east, color=bridgegray,
    font=\footnotesize\itshape\sffamily,
    inner sep=2pt
  },
  caption/.style={
    align=center, font=\scriptsize\itshape\sffamily,
    color=bridgegray, inner sep=1pt
  },
  flowarrow/.style={-{Stealth[length=3mm]}, thick, draw=bridgeblue}
]

% =====================================================================
% Etiquettes des entrees (a gauche)
% =====================================================================
\node[inputlbl] (lblTcp) at (-0.2, 4.5) {TCP / HL7v2};
\node[inputlbl] (lblBle) at (-0.2, 3.7) {BLE / MQTT};
\node[inputlbl] (lblFil) at (-0.2, 2.9) {Fichier XML};
\node[inputlbl] (lblRs)  at (-0.2, 2.1) {RS-232};

% =====================================================================
% Couche 1 : Adapter (teal pour distinguer la frontiere d'entree)
% =====================================================================
\node[altlayer] (adapter) at (1.4, 3.3) {Adapter\\Layer};

% Fleches d'entree
\draw[flowarrow] (lblTcp.east) -- (adapter.west |- lblTcp.east);
\draw[flowarrow] (lblBle.east) -- (adapter.west |- lblBle.east);
\draw[flowarrow] (lblFil.east) -- (adapter.west |- lblFil.east);
\draw[flowarrow] (lblRs.east)  -- (adapter.west |- lblRs.east);

% =====================================================================
% Couche 2 : Parser
% =====================================================================
\node[layer] (parser) at (4.6, 3.3) {Parser\\Layer};
\draw[flowarrow] (adapter.east) -- (parser.west)
  node[midway, above, font=\scriptsize\itshape, color=bridgegray, yshift=2pt]
    {trame brute};

% =====================================================================
% Couche 3 : Mapper
% =====================================================================
\node[layer] (mapper) at (7.8, 3.3) {Mapper\\Layer};
\draw[flowarrow] (parser.east) -- (mapper.west)
  node[midway, above, font=\scriptsize\itshape, color=bridgegray, yshift=2pt]
    {objet};

% =====================================================================
% Couche 4 : Transport
% =====================================================================
\node[altlayer] (transp) at (11.0, 3.3) {Transport\\Layer};
\draw[flowarrow] (mapper.east) -- (transp.west)
  node[midway, above, font=\scriptsize\itshape, color=bridgegray, yshift=2pt]
    {FHIR};

% =====================================================================
% Sortie HAPI
% =====================================================================
\node[outbox] (hapi) at (13.9, 3.3) {HAPI FHIR};
\draw[flowarrow] (transp.east) -- (hapi.west)
  node[midway, above, font=\scriptsize\itshape, color=bridgegray, yshift=2pt]
    {mTLS};

% =====================================================================
% Annotations en bas (bien sous chaque couche, espace suffisant)
% =====================================================================
\node[caption] at (1.4, 0.7) {drivers\\d'entree};
\node[caption] at (4.6, 0.7) {profils JSON\\decodage};
\node[caption] at (7.8, 0.7) {LOINC\\Observation FHIR};
\node[caption] at (11.0, 0.7) {retry, fallback\\file SQLite};

% =====================================================================
% Bandeau plugins (au-dessus, couvrant les 3 couches extensibles)
% =====================================================================
\node[draw=bridgegray, dashed, fill=bridgegray!8, rounded corners=2pt,
      minimum width=8.4cm, minimum height=0.7cm,
      anchor=south, font=\scriptsize\itshape\sffamily, color=bridgegray]
  at (4.6, 6.0)
  {Plugins communautaires : extensions Adapter, Parser, Mapper};

% Petites fleches verticales du bandeau vers les 3 couches
\foreach \x in {1.4, 4.6, 7.8} {
  \draw[->, dashed, color=bridgegray, thin] (\x, 5.95) -- (\x, 5.45);
}

\end{tikzpicture}
}
  \caption{Pipeline fonctionnel à 4 couches de la passerelle, avec annotations des formats à chaque interface.}
  \label{fig:rapport-pipeline}
\end{figure}

Chaque couche est un point d'extension. Un nouveau dispositif s'ajoute par un driver d'Adapter (canal nouveau), un module Parser (grammaire nouvelle) et un profil JSON. Transport et format de sortie restent stables. Cette séparation autorise une croissance par plugins (roadmap F1).

\section{Workflow d'intégration au paysage belge}
\label{sec:workflow-belge}

Le scénario complet mobilise quatre maillons : identification patient via NISS et services de base eHealth (simulé par identifiant de chambre en démo), soumission FHIR au DPI hospitalier (HAPI FHIR en démo), publication conditionnelle au MetaHub via le hub régional compétent, consultation patient via MaSanté.be.

\begin{bridgeinfo}[Hors-scope démo, mentionné par cohérence]
Le hub régional, eHealthBox et MaSanté.be ne sont pas implémentés dans la démonstration. Le pipeline démontré s'arrête à la soumission FHIR vers le serveur HAPI local. Le passage à l'échelle suppose un certificat eHealth de l'institution et une intégration au hub régional compétent.
\end{bridgeinfo}

\section{Comparaison synthétique des standards}
\label{sec:standards-synthese}

\begin{table}[h]
  \centering
  \caption{Standards mobilisés et leur position dans la passerelle.}
  \label{tab:standards-synthese}
  \bridgezebra
  \begin{tabular}{p{2.0cm} p{3.4cm} p{3.0cm} p{4.8cm}}
    \toprule
    \bridgeheader Standard & Domaine & Format & Position dans la passerelle \\
    \midrule
    HL7 v2 & Messages cliniques & ASCII pipe-delimited & Entrée Adapter (simulateur Philips) \\
    FHIR R4 & Ressources REST & JSON & Sortie principale du Mapper et du Transport \\
    IHE PCD & Profil d'intégration & Cadre de profils & Inspiration architecturale \\
    LOINC & Codes de mesures & Codes courts & Codification dans le Mapper \\
    KMEHR & Échange belge & XML & Hors-scope démo, roadmap \\
    DICOM & Imagerie médicale & Format binaire & Hors-scope démo, roadmap \\
    SNOMED-CT & Terminologie clinique & Codes longs & Hors-scope démo, roadmap \\
    \bottomrule
  \end{tabular}
\end{table}

% =====================================================================
% Chapitre 4 : Architecture technologique de l'e-health
% =====================================================================
\chapter{Architecture technologique de l'e-health}
\label{ch:reseau}

Ce chapitre relit le projet par le modèle OSI, justifie l'architecture edge, expose les compromis réseau et détaille la chaîne de sécurité.

\section{Le modèle OSI comme grille de lecture}
\label{sec:osi}

Le modèle OSI à sept couches situe chaque protocole et identifie la couche responsable d'un dysfonctionnement. La passerelle traverse les sept couches : trames hétérogènes en couches basses côté dispositif, messages standardisés en couches hautes côté plateforme.

\begin{table}[h]
  \centering
  \caption{Pile OSI traduite par la passerelle, du dispositif Legacy à la plateforme moderne.}
  \label{tab:osi-pile}
  \bridgezebra
  \begin{tabular}{p{3.0cm} p{5.5cm} p{5.5cm}}
    \toprule
    \bridgeheader Couche & Côté dispositif Legacy & Côté plateforme moderne \\
    \midrule
    7 Application & Format propriétaire ou HL7 v2 & HL7 FHIR R4 (REST) \\
    6 Présentation & ASCII binaire propriétaire & JSON ou OAuth 2.0 \\
    5 Session & SIP éventuel & HTTPS persistante \\
    4 Transport & Série ou TCP brut & TCP + TLS 1.3 \\
    3 Réseau & IPv4 local & IPv4 ou IPv6 routée \\
    2 Liaison & RS-232, RS-485, Bluetooth Classic & Ethernet ou Wi-Fi WPA3 \\
    1 Physique & DB-9, USB, ondes 2,4 GHz & Cat 6, fibre ou 5G \\
    \bottomrule
  \end{tabular}
\end{table}

La passerelle est un traducteur de pile : elle absorbe une pile pauvre, faiblement sécurisée et non routable, et émet sur une pile riche, chiffrée et fédérée. Cette lecture par couches structure aussi le diagnostic en exploitation.

\section{Chaîne canonique capteur-application-actionneur}
\label{sec:chaine-capteur}

Le cours TS6 propose une chaîne fonctionnelle universelle pour situer toute brique d'un système IoT médical. La passerelle s'y substitue au maillon \emph{application}, sans toucher aux maillons amont.

\begin{figure}[h]
  \centering
  \resizebox{\linewidth}{!}{% =====================================================================
% figures/fig-chaine-capteur-actionneur.tex
% Chaine canonique du cours : capteur -> microcontroleur -> communication
%                          -> application -> actionneur
% La passerelle se substitue au maillon Application.
% =====================================================================
\begin{tikzpicture}[
  font=\sffamily,
  every node/.style={font=\footnotesize\sffamily},
  >={Stealth[length=2.5mm]},
  thick,
  block/.style={
    rectangle, rounded corners=2pt,
    draw=bridgeblue, fill=bridgeblue!8, text=bridgeblue,
    minimum width=2.5cm, minimum height=1.1cm, align=center,
    inner sep=3pt, font=\footnotesize\bfseries\sffamily
  },
  highlight/.style={
    rectangle, rounded corners=2pt,
    draw=bridgeteal, fill=bridgeteal!15, text=bridgeteal,
    minimum width=2.5cm, minimum height=1.1cm, align=center,
    inner sep=3pt, font=\footnotesize\bfseries\sffamily
  },
  flow/.style={-{Stealth[length=3mm]}, thick, draw=bridgeblue}
]

% --- Chaine horizontale ----------------------------------------------
\node[block] (capt) at (0, 0) {Capteur};
\node[block] (mcu)  at (3.0, 0) {Micro-\\contr\^oleur};
\node[block] (comm) at (6.0, 0) {Module de\\communication};
\node[highlight] (app) at (9.0, 0) {Application\\(notre passerelle)};
\node[block] (act)  at (12.0, 0) {Actionneur};

\draw[flow] (capt.east) -- (mcu.west);
\draw[flow] (mcu.east)  -- (comm.west);
\draw[flow] (comm.east) -- (app.west);
\draw[flow] (app.east)  -- (act.west);

% --- Annotations sous les blocs --------------------------------------
\node[font=\scriptsize\itshape, color=bridgegray, align=center]
  at ([yshift=-0.85cm]capt.south) {grandeur\\physique};
\node[font=\scriptsize\itshape, color=bridgegray, align=center]
  at ([yshift=-0.85cm]mcu.south) {numerise\\conditionne};
\node[font=\scriptsize\itshape, color=bridgegray, align=center]
  at ([yshift=-0.85cm]comm.south) {RS-232, BLE,\\TCP, fichier};
\node[font=\scriptsize\itshape, color=bridgeteal, align=center]
  at ([yshift=-0.85cm]app.south) {capture passive,\\transformation FHIR};
\node[font=\scriptsize\itshape, color=bridgegray, align=center]
  at ([yshift=-0.85cm]act.south) {pompe,\\alarme, etc.};

% --- Surbrillance "passerelle" ---------------------------------------
\draw[bridgeteal, thick, dashed, rounded corners=4pt]
  ([xshift=-0.3cm, yshift=0.4cm]app.north west) rectangle
  ([xshift=0.3cm, yshift=-0.4cm]app.south east);
\node[anchor=south, font=\scriptsize\bfseries, color=bridgeteal]
  at ([yshift=0.45cm]app.north) {Substitution};

\end{tikzpicture}
}
  \caption{Chaîne canonique IoT médicale et position de la passerelle.}
  \label{fig:rapport-chaine-capteur}
\end{figure}

Côté réseau, le parcours d'un message traverse plusieurs équipements actifs : commutateurs, routeurs, points d'accès Wi-Fi, et la passerelle elle-même au sens applicatif. La sécurité de bout en bout repose donc sur la coordination de ces éléments, et non sur un seul nœud.

\section{Mesures réseau et compromis}
\label{sec:mesures}

Trois mesures décrivent le transport : latence (ms), débit (octets/s) et bande passante. Ces grandeurs entrent en compromis selon le cas d'usage.

\begin{table}[h]
  \centering
  \caption{Cas d'usage et priorités réseau associées.}
  \label{tab:cas-usage}
  \bridgezebra
  \begin{tabular}{p{6.0cm} p{8.0cm}}
    \toprule
    \bridgeheader Cas d'usage & Priorités \\
    \midrule
    Télémonitoring chronique (notre cas) & Fiabilité, autonomie, débit modéré \\
    Bloc opératoire connecté & Latence basse, déterminisme strict \\
    Échographe portable & Bande passante élevée \\
    Glucomètre point of care & Autonomie batterie, fiabilité ponctuelle \\
    \bottomrule
  \end{tabular}
\end{table}

\section{Architecture edge plutôt que cloud}
\label{sec:edge}

Quatre arguments justifient le bord de réseau hospitalier : latence locale sous-milliseconde côté dispositif, souveraineté des données (flux jamais sortis du périmètre sans contrôle), robustesse en cas de coupure WAN (buffer local persistant), respect du principe RGPD de minimisation. Le choix d'un middleware physique répond aux contraintes matérielles : compatibilité ports série historiques, isolation des dispositifs critiques sur VLAN dédié, surface d'attaque réduite (la passerelle est le seul équipement de la zone à parler vers l'extérieur).

\section{Sécurité réseau}
\label{sec:secu-reseau}

La sécurité du transport s'articule autour de quatre dispositifs complémentaires.

\subsection{Segmentation et VLAN}

Les dispositifs sont placés sur un VLAN dédié, isolé du reste du réseau hospitalier. La passerelle est le seul point de sortie autorisé. Cette segmentation respecte le principe de cloisonnement de l'ISO 27799 \cite{iso_27799} et limite le mouvement latéral en cas de compromission.

\subsection{Chiffrement en transit}

TLS 1.3 en sortie apporte confidentialité, intégrité et résistance aux attaques par rejeu. mTLS est imposé vers le DPI, conformément à l'ISO 13608 \cite{iso_13608}. Les liaisons Legacy en entrée (RS-232, BLE classique) ne sont pas chiffrables nativement : VLAN dédié et capture passive constituent alors la défense.

\subsection{Services ETEE de la plateforme eHealth}

Pour les échanges via MetaHub ou eHealthBox, la plateforme eHealth fournit les services ETKDepot (dépôt de clés publiques) et KGSS (génération de clés de session) pour chiffrer un message KMEHR \cite{ehealth_services_base}. Hors-scope démo, mentionné pour la mise en production.

\subsection{Authentification de l'institution}

L'institution dispose d'un certificat eHealth comme identité forte vers BIHR, MetaHub et eHealthBox. Pour les échanges privés vers un DPI propriétaire, mTLS sortant s'applique avec un certificat émis par l'AC interne de l'hôpital.

\section{Pipeline réseau de la passerelle}
\label{sec:pipeline-reseau}

La chaîne Adapter, Parser, Mapper, Transport (cf. § 3.7) produit ou consomme un format différent à chaque couche, et la sécurité s'applique uniformément en aval du Mapper. Le pipeline complet, files de repli locales et rejeu en cas d'échec, est documenté en annexe B du cahier des charges.

\begin{figure}[h]
  \centering
  \resizebox{\linewidth}{!}{% =====================================================================
% figures/fig-pipeline-4couches.tex
% Pipeline fonctionnel : Adapter -> Parser -> Mapper -> Transport -> HAPI
% Inputs varies a gauche, sortie FHIR a droite. Bandeau plugins au-dessus.
% =====================================================================
\begin{tikzpicture}[
  font=\sffamily,
  every node/.style={font=\small\sffamily},
  >={Stealth[length=2.5mm]},
  thick,
  layer/.style={
    rectangle, rounded corners=3pt,
    draw=bridgeblue, fill=bridgeblue!10, text=bridgeblue,
    minimum width=2.4cm, minimum height=4.2cm,
    align=center,
    font=\small\bfseries\sffamily
  },
  altlayer/.style={
    rectangle, rounded corners=3pt,
    draw=bridgeteal, fill=bridgeteal!10, text=bridgeteal,
    minimum width=2.4cm, minimum height=4.2cm,
    align=center,
    font=\small\bfseries\sffamily
  },
  outbox/.style={
    rectangle, rounded corners=3pt,
    draw=bridgegreen, fill=bridgegreen!15, text=bridgegreen,
    minimum width=2.0cm, minimum height=1.0cm,
    align=center,
    font=\small\bfseries\sffamily
  },
  inputlbl/.style={
    anchor=east, color=bridgegray,
    font=\footnotesize\itshape\sffamily,
    inner sep=2pt
  },
  caption/.style={
    align=center, font=\scriptsize\itshape\sffamily,
    color=bridgegray, inner sep=1pt
  },
  flowarrow/.style={-{Stealth[length=3mm]}, thick, draw=bridgeblue}
]

% =====================================================================
% Etiquettes des entrees (a gauche)
% =====================================================================
\node[inputlbl] (lblTcp) at (-0.2, 4.5) {TCP / HL7v2};
\node[inputlbl] (lblBle) at (-0.2, 3.7) {BLE / MQTT};
\node[inputlbl] (lblFil) at (-0.2, 2.9) {Fichier XML};
\node[inputlbl] (lblRs)  at (-0.2, 2.1) {RS-232};

% =====================================================================
% Couche 1 : Adapter (teal pour distinguer la frontiere d'entree)
% =====================================================================
\node[altlayer] (adapter) at (1.4, 3.3) {Adapter\\Layer};

% Fleches d'entree
\draw[flowarrow] (lblTcp.east) -- (adapter.west |- lblTcp.east);
\draw[flowarrow] (lblBle.east) -- (adapter.west |- lblBle.east);
\draw[flowarrow] (lblFil.east) -- (adapter.west |- lblFil.east);
\draw[flowarrow] (lblRs.east)  -- (adapter.west |- lblRs.east);

% =====================================================================
% Couche 2 : Parser
% =====================================================================
\node[layer] (parser) at (4.6, 3.3) {Parser\\Layer};
\draw[flowarrow] (adapter.east) -- (parser.west)
  node[midway, above, font=\scriptsize\itshape, color=bridgegray, yshift=2pt]
    {trame brute};

% =====================================================================
% Couche 3 : Mapper
% =====================================================================
\node[layer] (mapper) at (7.8, 3.3) {Mapper\\Layer};
\draw[flowarrow] (parser.east) -- (mapper.west)
  node[midway, above, font=\scriptsize\itshape, color=bridgegray, yshift=2pt]
    {objet};

% =====================================================================
% Couche 4 : Transport
% =====================================================================
\node[altlayer] (transp) at (11.0, 3.3) {Transport\\Layer};
\draw[flowarrow] (mapper.east) -- (transp.west)
  node[midway, above, font=\scriptsize\itshape, color=bridgegray, yshift=2pt]
    {FHIR};

% =====================================================================
% Sortie HAPI
% =====================================================================
\node[outbox] (hapi) at (13.9, 3.3) {HAPI FHIR};
\draw[flowarrow] (transp.east) -- (hapi.west)
  node[midway, above, font=\scriptsize\itshape, color=bridgegray, yshift=2pt]
    {mTLS};

% =====================================================================
% Annotations en bas (bien sous chaque couche, espace suffisant)
% =====================================================================
\node[caption] at (1.4, 0.7) {drivers\\d'entree};
\node[caption] at (4.6, 0.7) {profils JSON\\decodage};
\node[caption] at (7.8, 0.7) {LOINC\\Observation FHIR};
\node[caption] at (11.0, 0.7) {retry, fallback\\file SQLite};

% =====================================================================
% Bandeau plugins (au-dessus, couvrant les 3 couches extensibles)
% =====================================================================
\node[draw=bridgegray, dashed, fill=bridgegray!8, rounded corners=2pt,
      minimum width=8.4cm, minimum height=0.7cm,
      anchor=south, font=\scriptsize\itshape\sffamily, color=bridgegray]
  at (4.6, 6.0)
  {Plugins communautaires : extensions Adapter, Parser, Mapper};

% Petites fleches verticales du bandeau vers les 3 couches
\foreach \x in {1.4, 4.6, 7.8} {
  \draw[->, dashed, color=bridgegray, thin] (\x, 5.95) -- (\x, 5.45);
}

\end{tikzpicture}
}
  \caption{Pipeline fonctionnel de la passerelle, lu sous l'angle réseau.}
  \label{fig:rapport-pipeline-reseau}
\end{figure}

\section{Perspective : SDN et orchestration multi-gateways}
\label{sec:sdn}

Un grand hôpital peut compter plusieurs dizaines de passerelles par étage ou service. Une approche SDN ou une orchestration Kubernetes permet le provisioning centralisé, les mises à jour et la micro-segmentation dynamique. Évolution traitée au chapitre 14 du cahier des charges.

% =====================================================================
% Chapitre 5 : Capteurs et dispositifs de mesure en e-health
% =====================================================================
\chapter{Capteurs et dispositifs de mesure en e-health}
\label{ch:capteurs}

Ce chapitre typologise les capteurs par grandeur physique, présente les quatre équipements de référence et détaille les méthodes de capture passive qui préservent le marquage CE.

\section{Définition canonique du capteur}
\label{sec:def-capteur}

Un capteur transforme une grandeur physique en signal électrique exploitable. La chaîne de mesure type comporte cinq maillons : capteur, microcontrôleur (conditionnement, numérisation), module de communication, application, actionneur. La passerelle se branche au point de communication, jamais au capteur. Elle n'altère ni le conditionnement analogique ni le traitement numérique embarqué, ce qui préserve la conformité MDR du dispositif source.

\section{Typologie par grandeur physique}
\label{sec:typologie}

\begin{table}[h]
  \centering
  \caption{Familles de capteurs, grandeurs mesurées et exemples du projet.}
  \label{tab:typologie-capteurs}
  \bridgezebra
  \begin{tabular}{p{3.5cm} p{5.0cm} p{6.0cm}}
    \toprule
    \bridgeheader Famille & Grandeur mesurée & Exemples projet \\
    \midrule
    Optique & Saturation O2 par photopléthysmographie & Oxymètre Masimo Radical-7 \\
    Électrique (biopotentiel) & ECG, fréquence cardiaque & Moniteur Philips IntelliVue, ECG Welch Allyn CP150 \\
    Pression & Pression artérielle oscillométrique, PEEP & Moniteur Philips, ventilateur Dräger \\
    Mécanique & Volumes respiratoires & Spiromètre Vyaire MicroLab \\
    Chimique (point of care) & Glycémie, gaz du sang & Glucomètre Roche Accu-Chek \\
    Température & Température cutanée ou interne & Thermomètres connectés (perspective) \\
    \bottomrule
  \end{tabular}
\end{table}

\section{Équipements de référence retenus pour la démonstration}
\label{sec:equipements-demo}

Quatre équipements ont été sélectionnés, représentatifs des canaux d'entrée et bénéficiant d'une documentation publique solide. Détails techniques dans \texttt{EQUIPMENTS.md}.

\begin{table}[h]
  \centering
  \caption{Équipements de référence et cibles LOINC associées.}
  \label{tab:equipements-demo}
  \bridgezebra
  \begin{tabular}{p{3.5cm} p{3.0cm} p{3.5cm} p{4.0cm}}
    \toprule
    \bridgeheader Équipement & Famille & Sortie native & Cible LOINC \\
    \midrule
    Philips IntelliVue MX450 & Moniteur multiparamétrique & HL7 v2 ORU\textasciicircum{}R01 over MLLP & SpO2 (59408-5), FC (8867-4), PA (8480-6, 8462-4) \\
    Masimo Radical-7 & Oxymètre point of care & JSON via BLE simulé MQTT & SpO2 (59408-5), FC (8867-4) \\
    Welch Allyn CP150 & ECG 12 dérivations & Fichier SCP-ECG ISO 11073-91064 & Étude ECG (11524-6) \\
    Dräger Evita V500 & Ventilateur & ASCII MEDIBUS sur RS-232 & FiO2 (19994-3), PEEP (20077-4) \\
    \bottomrule
  \end{tabular}
\end{table}

Documentation constructeur exploitable : Philips IntelliVue MX450 et IntelliVue Information Center iX \cite{philips_piic_ix}, Masimo Radical-7 manuel opérateur \cite{masimo_radical7_manual}, Welch Allyn CP150 notice et manuel de service \cite{welchallyn_cp150_dfu}, Dräger Evita V500 protocole MEDIBUS (ASCII, trame STX/ETX) \cite{draeger_medibus_protocol}.

\section{Méthodes de capture sans intrusion matérielle}
\label{sec:capture-non-intrusive}

\begin{table}[h]
  \centering
  \caption{Méthodes de capture passive employées dans la démonstration.}
  \label{tab:capture-non-intrusive}
  \bridgezebra
  \begin{tabular}{p{4.0cm} p{10.5cm}}
    \toprule
    \bridgeheader Méthode & Description \\
    \midrule
    Capture passive sur port série & Lecture seule en pseudo-tty ou null-modem. Aucune modification du dispositif. Aucune perte du marquage CE. Notre cas pour Dräger. \\
    Tap réseau passif & SPAN ou TAP sur le commutateur. Lecture des trames TCP émises par le dispositif. Notre cas pour Philips et Masimo. \\
    Lecture de fichiers exportés & Surveillance d'un dossier où le dispositif dépose ses CSV ou XML. Notre cas pour Welch Allyn. \\
    OCR sur écran & Capture d'écran traitée par un modèle de reconnaissance. Recours ultime, fiabilité limitée. Hors-scope démonstration. \\
    \bottomrule
  \end{tabular}
\end{table}

\begin{bridgekey}[Principe de non-intrusion]
Aucune méthode retenue ne modifie le dispositif. Le marquage CE reste valide. Le dispositif ne sait pas qu'il est observé. Cette posture est le fondement réglementaire de la passerelle au regard du règlement MDR \cite{mdr_2017_745}.
\end{bridgekey}

\section{Capteurs émergents et perspectives IA}
\label{sec:capteurs-emergents}

De nouveaux capteurs intègrent une couche d'analyse en bord : analyse de frappe clavier (détection précoce de troubles cognitifs), vision par ordinateur sur caméra de chambre (détection automatique de chute). Notre roadmap F2 prévoit une couche d'anomaly detection en bord qui pré-flagge les valeurs aberrantes avant transmission.

\section{Boucle vers le DPI}
\label{sec:boucle-dpi}

Tout résultat issu d'un dispositif médical remonte au DPI. Une fois mappée en FHIR R4 et publiée vers le serveur cible, la donnée est consultée par le clinicien via le DPI ou par le patient via MaSanté. La passerelle disparaît derrière la valeur clinique livrée, effet attendu d'un middleware réussi.

\chapter{Implémentation, adoption, évaluation et perspectives}
\label{ch:deploiement}

Comment déployer la passerelle, en faire adopter l'usage et prouver sa valeur ? Ce chapitre articule calendrier de déploiement, statut MDR, cadre mHealth Belgium, sept facteurs d'adoption et grille d'évaluation. Il prépare la dimension sécurité du chapitre~\ref{ch:ethique-securite}.

\section{Phases d'implémentation}
\label{sec:deploiement-phases}

\bridgezebra
\begin{tabular}{p{2.2cm}p{6.5cm}p{5cm}}
\toprule
\bridgeheader Phase & Description & Livrable \\
\midrule
Phase 0 & Audit du parc dispositifs hospitaliers, identification des familles cibles & Cartographie OT \\
Phase 1 & POC sur un service (USI ou pneumologie), un seul type de dispositif & Démo dockerisée (présent travail) \\
Phase 2 & Extension à 3 à 5 familles, intégration au DPI test & Pilote unitaire \\
Phase 3 & Pilote multisite, raccordement au MetaHub & Pilote intégré \\
Phase 4 & Industrialisation, passage en production & Déploiement opérationnel \\
\bottomrule
\end{tabular}

La démonstration livrée correspond à la phase~1. Les phases ultérieures relèvent d'un projet d'industrialisation porté conjointement par la DSI, le service biomédical et le DPO.

\section{Statut réglementaire de la passerelle}
\label{sec:deploiement-mdr}

\begin{bridgekey}[Statut non-DM revendiqué]
La passerelle est un middleware d'intégration IT. Elle ne décide rien cliniquement. Elle ne modifie pas les dispositifs amont. Elle consomme des sorties existantes et les normalise. Cette posture la place hors du champ MDR pour la passerelle elle-même, tout en respectant le MDR pour les dispositifs Legacy connectés grâce à une capture passive sans modification firmware.
\end{bridgekey}

Conséquences : aucun marquage CE pour le middleware, aucune notification AFMPS au titre du middleware. La conformité MDR reste portée par les dispositifs amont. La conformité du DPI aval relève du responsable de traitement hospitalier \cite{mdr_2017_745,afmps}.

\section{Pyramide mHealth Belgium adaptée}
\label{sec:deploiement-mhealth}

Le cadre mhealthbelgium.be propose trois niveaux M1, M2, M3 pour les apps de santé. Bien que conçu pour les apps, il fournit une grille de lecture utile au middleware \cite{mhealthbelgium_pyramide}.

\bridgezebra
\begin{tabular}{p{1.2cm}p{5cm}p{7.8cm}}
\toprule
\bridgeheader Niveau & Critères & Application au projet \\
\midrule
M1 & Exigences minimales (CE, AFMPS, RGPD) & Capture passive MDR-friendly, conformité RGPD, pas de marquage CE requis (statut non-DM) \\
M2 & Sécurité, authentification, interopérabilité & mTLS, certificats eHealth, FHIR R4, profil IHE PCD \\
M3 & Remboursement, valeur ajoutée prouvée & Hors-scope démo, dépend de l'industrialisation \\
\bottomrule
\end{tabular}

La passerelle satisfait par construction M1 et vise M2 dès la mise en production. M3 dépend d'études médico-économiques hors stade démonstrateur.

\section{Sept facteurs d'adoption appliqués}
\label{sec:deploiement-adoption}

\bridgezebra
\begin{tabular}{p{4.5cm}p{9.5cm}}
\toprule
\bridgeheader Facteur & Mise en oeuvre projet \\
\midrule
Compréhension et formation & Documentation pour techniciens biomédicaux, dashboard explicite \\
Interopérabilité & Coeur du projet : FHIR R4, LOINC, profil IHE PCD \\
Engagement parties prenantes & Co-conception avec biomédicaux, médecins, infirmiers, DPO en phase prod \\
Accessibilité et UX & Dashboard sobre, statut couleur immédiat, alertes mail \\
Support et maintenance & Mises à jour OTA, profils JSON modifiables sans redémarrage \\
Conformité & MDR (capture passive), RGPD (minimisation), certifications eHealth (en prod) \\
Intégration aux flux & Pas de saisie manuelle ajoutée, données arrivent automatiquement au DPI \\
\bottomrule
\end{tabular}

\section{Évaluation et indicateurs}
\label{sec:deploiement-evaluation}

Quatre catégories d'indicateurs structurent la preuve de valeur, chacune adressant une partie prenante (DSI, soignants, médecins, direction financière).

\bridgezebra
\begin{tabular}{p{3cm}p{6cm}p{5cm}}
\toprule
\bridgeheader Catégorie & Indicateur & Objectif \\
\midrule
Technique & Taux de transmission réussie & supérieur à 99\% \\
Technique & Latence p95 bout-en-bout & inférieure à 500 ms \\
Technique & Taux d'erreur de mapping & inférieur à 1\% \\
Organisationnel & Minutes par jour économisées sur la saisie manuelle & supérieur à 30 min par soignant \\
Organisationnel & Taux de re-saisie évité & supérieur à 90\% \\
Clinique & Complétude du dossier patient & hausse mesurable trimestrielle \\
Médico-économique & Coût par lit comparé à un parc neuf nativement connecté & inférieur d'un facteur 5 \\
\bottomrule
\end{tabular}

L'argument médico-économique est central : la passerelle évite de jeter quinze ans d'investissement matériel. Le retour sur investissement se mesure en mois.

\section{Perspectives industrialisation}
\label{sec:deploiement-industrialisation}

Le passage à l'échelle suppose un déploiement multi-services et multi-sites avec fédération de gateways supervisées centralement (cf. chapitre 14 du cahier des charges). Trois axes : orchestration SDN pour gérer un parc, plugins communautaires pour étendre le catalogue, fédération européenne EHDS pour le partage transfrontalier.

\chapter{Éthique, sécurité et confidentialité}
\label{ch:ethique-securite}

Comment garantir confidentialité, intégrité, disponibilité, traçabilité et conformité légale tout au long du cycle de vie des données ? Ce chapitre articule la triade CIA augmentée de la traçabilité, la conformité RGPD, le cycle de vie de la donnée, la cartographie des risques, les services ETEE, la responsabilité éthique, la gouvernance et les normes ISO.

\section{Triade CIA et traçabilité}
\label{sec:ethique-cia}

La sécurité de l'information en santé repose sur trois piliers historiques (confidentialité, intégrité, disponibilité), augmentés de la traçabilité dès lors que la responsabilité clinique est engagée.

\bridgezebra
\begin{tabular}{p{3cm}p{5cm}p{6.5cm}}
\toprule
\bridgeheader Pilier & Définition & Mise en oeuvre projet \\
\midrule
Confidentialité & Accès restreint aux personnes autorisées & Chiffrement TLS, principe de minimisation, dashboard sans donnée patient \\
Intégrité & Données non altérées & Checksum SHA-256 sur file de repli, validation FHIR avant émission \\
Disponibilité & Système accessible quand requis & Retry exponentiel, fallback SQLite, hot-reload des profils \\
Traçabilité & Chaque action est journalisée & Logs JSON horodatés UTC, sans valeur clinique \\
\bottomrule
\end{tabular}

\section{Conformité RGPD}
\label{sec:ethique-rgpd}

Le RGPD encadre tout traitement de données personnelles, et particulièrement les données de santé (catégorie particulière, art. 9).

\bridgezebra
\begin{tabular}{p{5.5cm}p{8.5cm}}
\toprule
\bridgeheader Obligation RGPD & Mise en oeuvre \\
\midrule
Registre des traitements (art. 30) & Documenté côté hôpital, esquissé en annexe du CDC \\
Base légale (art. 6.1.e et 9.2.h) & Mission d'intérêt public en santé, exécution d'un contrat de soin \\
Minimisation (art. 5) & Aucune donnée patient dans logs, dashboard, alertes \\
Durée de conservation (art. 5) & Déterminée par le responsable de traitement (l'hôpital) \\
Notification de violation (art. 33) & Procédure 72h documentée au CDC chapitre 12 \\
Analyse d'impact (art. 35) & AIPD esquissée au CDC chapitre 12, AIPD complète obligatoire en prod \\
Information des personnes (art. 13-14) & Portée par l'hôpital via la politique générale du DPI \\
\bottomrule
\end{tabular}

Sources : \cite{rgpd_2016_679}, \cite{loi_protection_donnees_2018}, \cite{loi_droits_patient_2002}.

Le partage des données vers les hubs régionaux (RSW, ABRUMET, CoZo) est conditionné à un consentement explicite du patient au partage hub, distinct du consentement RGPD général au traitement \cite{esante_wallonie_rgpd}.

\section{Cycle de vie de la donnée de santé}
\label{sec:ethique-cycle}

Le cycle de vie comporte cinq étapes (écriture, lecture, modification, conservation, destruction), chacune portant ses risques et mesures.

\bridgezebra
\begin{tabular}{p{3.5cm}p{10.5cm}}
\toprule
\bridgeheader Étape & Mesure projet \\
\midrule
Écriture (capture) & Capture passive sans modification du dispositif \\
Lecture et accès & Matrice d'accès du DPI selon le rôle, la relation thérapeutique vérifiée et le consentement du patient. \\
Modification & Impossible côté passerelle, qui est read-only \\
Conservation & File de repli locale chiffrée Fernet, durée limitée (24h max en démo) \\
Destruction & Suppression effective du fichier SQLite à la fin de la démo \\
\bottomrule
\end{tabular}

\section{Cartographie des risques}
\label{sec:ethique-risques}

Une analyse STRIDE simplifiée recense les risques majeurs (version détaillée au § 12 du CDC).

\bridgezebra
\begin{tabular}{p{3cm}p{5cm}p{6cm}}
\toprule
\bridgeheader Surface & Risque & Mesure \\
\midrule
Liaison Adapter & Injection de fausses observations & Validation Device ID, plages de plausibilité \\
Parser Layer & Buffer overflow ou DoS & Timeout, taille max, file de quarantaine \\
Mapper Layer & Mapping erroné, mauvaise unité & Validation FHIR, tests unitaires \\
Transport Layer & MITM sur HAPI & mTLS optionnel en démo, obligatoire en prod \\
Stockage local & Lecture du SQLite par un attaquant local & Chiffrement Fernet, OS minimal \\
Plugin tiers & Plugin malveillant & Validation manifeste TOML, sandbox d'exécution \\
\bottomrule
\end{tabular}

\section{Services ETEE de la plateforme eHealth}
\label{sec:ethique-etee}

Les services ETEE (ETKDepot, KGSS) assurent un chiffrement bout-en-bout entre acteurs autorisés. Hors-scope démo locale (qui s'arrête à mTLS vers HAPI), indispensables en production pour les hubs régionaux et MetaHub \cite{ehealth_services_base}.

\section{Question éthique : qualité des données et responsabilité}
\label{sec:ethique-responsabilite}

Si le mapping introduit une erreur (mauvaise unité, mauvais code LOINC), qui est responsable ? La chaîne se contractualise. La passerelle prend la responsabilité du mapping syntaxique et sémantique, validé par tests automatisés. Le dispositif prend la responsabilité de la mesure (marquage CE). L'hôpital prend la responsabilité de l'interprétation et de la décision clinique.

Un second risque éthique est la \emph{datafication} (collecter plus que nécessaire). La minimisation y répond : la passerelle ne capture que ce que le profil JSON déclare comme pertinent. Aucun flux n'est aspiré sans déclaration explicite.

\section{Gouvernance des données}
\label{sec:ethique-gouvernance}

Les rôles RGPD : responsable de traitement = hôpital qui déploie ; éditeur de la passerelle = sous-traitant (art. 28), lié par contrat de sous-traitance ; DPO impliqué dès la conception (privacy by design, art. 25) \cite{esante_wallonie_rgpd}.

\section{Normes ISO de référence}
\label{sec:ethique-iso}

Trois normes encadrent la sécurité de l'information en santé : ISO~13608 (sécurité des communications), ISO~27001 (SMSI), ISO~27799 (santé). Visées comme cibles de certification en production, non implémentées au stade démonstrateur \cite{iso_13608,iso_27001,iso_27799}.

% =====================================================================
% Conclusion : synthese a trois niveaux (technique, systemique, prospectif)
% =====================================================================
\chapter*{Conclusion}
\addcontentsline{toc}{chapter}{Conclusion}
\label{ch:conclusion}

Ce travail articule grille de lecture théorique (sept chapitres du cours TS6) et démonstration pratique (passerelle Edge dockerisée transformant quatre flux dispositifs en FHIR R4). La synthèse s'articule en trois niveaux.

\section*{Niveau technique}

Une passerelle Edge multi-protocoles transforme un parc Legacy en source FHIR R4 native, sans modification matérielle. L'architecture en quatre couches (Adapter, Parser, Mapper, Transport) est chargée dynamiquement à partir de profils JSON. Chaque couche est un point d'extension : ajouter un dispositif consiste à publier un profil déclaratif et, si besoin, un parseur Python en plugin. Cette séparation déclarative entre canal, langue et mapping est la traduction directe des concepts d'interopérabilité technique, syntaxique et sémantique du chapitre~\ref{ch:standards}.

La démonstration prouve la faisabilité sur quatre canaux : TCP HL7 v2 over MLLP (moniteurs IP), MQTT (BLE relayés), SCP-ECG déposés (ambulatoire), RS-232 MEDIBUS (Dräger). HAPI FHIR R4 reçoit des Bundles transactionnels conformes. Le dashboard supervise les flux sans exposer de donnée patient. Le mTLS sortant et la segmentation réseau (chapitre~\ref{ch:reseau}) bouclent la chaîne de confiance.

\section*{Niveau systémique}

La solution adresse un goulet d'étranglement national identifié par le Plan eSanté 2025-2027 : l'intégration automatique des dispositifs Legacy au DPI et au BIHR \cite{plan_esante_2025_2027,bihr_overview}. Les paramètres vitaux remontent encore par re-saisie ou exports manuels. Le BIHR ne peut fédérer que ce que les hôpitaux poussent. La passerelle ferme la boucle entre la mesure et le partage régional, et accélère le passage aux stades supérieurs de maturité numérique.

L'argument médico-économique est central : éviter de remplacer un parc valorisé sur 15 à 20 ans alors que les standards évoluent tous les 3 à 5 ans. La capture passive préserve le marquage CE conformément à l'esprit du Règlement (UE) 2017/745 \cite{mdr_2017_745}. Le statut non-DM revendiqué simplifie la chaîne de responsabilité : hôpital responsable de traitement, éditeur sous-traitant RGPD, fabricant amont responsable clinique du dispositif.

\section*{Niveau prospectif}

Trois évolutions dessinent l'avenir. F1 ouvre la passerelle à une communauté de plugins de traduction sur le modèle HACS (Home Assistant) et npm. Un fabricant ou hôpital qui rencontre un dispositif non couvert publie son plugin. F2 intègre une détection d'anomalies en bord : modèle léger entraîné hors-ligne, inférence ONNX au fil de l'eau, pré-flagging des valeurs aberrantes. F3 projette la solution dans l'European Health Data Space pour la mobilité transfrontalière (cf. chapitre 14 du cahier des charges).

L'effet de catalogue communautaire est le levier principal : aucun éditeur seul ne peut couvrir la diversité mondiale des dispositifs médicaux. Les conditions d'échelle sont une gouvernance ouverte, des mappings LOINC validés par revue de pairs, et des plugins signés par le RSSI hospitalier. La passerelle joue alors le rôle d'agrégateur de profils, à la manière des distributions Linux.

\section*{Mot final}

Le travail confirme que l'interopérabilité en santé n'est pas un obstacle technique mais une décision d'architecture. Le code, le cahier des charges et la démonstration sont remis au jury comme un ensemble cohérent.

% =====================================================================
% Annexes du rapport : figures, tableaux, renvoi CDC, glossaire, source demo
% =====================================================================
\appendix
\chapter{Annexes}
\label{ch:annexes-rapport}

\section*{Annexe A : Renvoi vers le cahier des charges}
\addcontentsline{toc}{section}{Annexe A : Renvoi vers le cahier des charges}

Le présent rapport décrit la démarche académique et la grille d'analyse fondée sur les sept chapitres du cours TS6. Le cahier des charges (\texttt{cdc.pdf}) précise les exigences techniques pour l'implémentation de la démonstration et trace la traçabilité entre exigences, profils JSON et tests de recette. Les deux documents se complètent et se lisent ensemble.

Les sections clés du CDC à consulter pour approfondir un point particulier sont rappelées dans le tableau~\ref{tab:renvoi-cdc}.

\begin{table}[!htbp]
\centering
\bridgezebra
\begin{tabular}{p{0.55\textwidth} p{0.35\textwidth}}
\toprule
\bridgeheader Sujet & Référence CDC \\
\midrule
Liste exhaustive des exigences fonctionnelles & CDC chapitre 4 (33 EF) \\
Exigences non fonctionnelles mesurables & CDC chapitre 5 (12 ENF) \\
Architecture détaillée et profils JSON & CDC chapitre 7 \\
Plan de réalisation et Gantt & CDC chapitre 9 \\
Recette finale de la démonstration & CDC chapitre 10 \\
AIPD esquissée et analyse STRIDE & CDC chapitre 12 \\
Roadmap F1 plugins, F2 anomalies, F3 EHDS & CDC chapitre 14 \\
Mapping LOINC complet et exemple Bundle FHIR & CDC chapitre 17 \\
\bottomrule
\end{tabular}
\caption{Renvois vers les sections du cahier des charges.}
\label{tab:renvoi-cdc}
\end{table}

\section*{Annexe B : Glossaire des sigles}
\addcontentsline{toc}{section}{Annexe B : Glossaire des sigles}

Le glossaire reprend les sigles employés dans le corps du rapport. Il complète le glossaire étendu du cahier des charges.

\begin{table}[!htbp]
\centering
\bridgezebra
\begin{tabular}{p{0.18\textwidth} p{0.72\textwidth}}
\toprule
\bridgeheader Sigle & Signification \\
\midrule
AFMPS & Agence Fédérale des Médicaments et Produits de Santé \\
AIPD & Analyse d'Impact relative à la Protection des Données \\
APD & Autorité de Protection des Données \\
BIHR & Belgian Integrated Health Record \\
BLE & Bluetooth Low Energy \\
CIA & Confidentialité, Intégrité, Disponibilité \\
DPI & Dossier Patient Informatisé \\
DPO & Data Protection Officer \\
EHDS & European Health Data Space \\
ETEE & End-to-End Encryption (services eHealth) \\
FHIR & Fast Healthcare Interoperability Resources \\
HL7 & Health Level Seven \\
IHE & Integrating the Healthcare Enterprise \\
INAMI & Institut National d'Assurance Maladie-Invalidité \\
INS & Identifiant National de Santé \\
KMEHR & Kind Messages for Electronic Healthcare Record \\
LOINC & Logical Observation Identifiers Names and Codes \\
MDR & Medical Device Regulation (UE 2017/745) \\
MEDIBUS & Protocole série Dräger pour ventilateurs et stations \\
MLLP & Minimal Lower Layer Protocol \\
MoSCoW & Méthode de priorisation Must, Should, Could, Won't \\
MQTT & Message Queuing Telemetry Transport \\
mTLS & Mutual Transport Layer Security \\
OT & Operational Technology \\
PCD & Patient Care Device (profil IHE) \\
RGPD & Règlement Général sur la Protection des Données \\
SCP-ECG & Standard Communications Protocol for ECG \\
TOML & Tom's Obvious Minimal Language \\
\bottomrule
\end{tabular}
\caption{Glossaire des sigles employés dans le rapport.}
\label{tab:glossaire-rapport}
\end{table}

\section*{Annexe C : Source de la démonstration}
\addcontentsline{toc}{section}{Annexe C : Source de la démonstration}

Le code source du middleware Python, les quatre simulateurs Docker, les profils JSON et le plugin exemple sont livrés sous forme d'un dépôt Git remis avec le présent rapport. La procédure de lancement tient en moins de dix lignes de commandes, documentées dans le \texttt{README.md} du dépôt. Le jury peut reproduire la démonstration en local avec \texttt{docker compose up}.


% --- Bibliographie ---------------------------------------------------
\printbibliography

\end{document}
